
Anyone setting out to prove a compiler correct must choose a style of
semantics.  The CakeML and CompCert projects, among others, show that
operational semantics such as abstract machines and interpreters are
viable choices. Their correctness proofs reason about program
execution using simulations and induction on evaluation. In this paper
we explore whether denotational semantics can also be used in proofs
of compiler correctness and evaluate the tradeoffs. However, a
literature review reveals an obstacle. Compiling a higher-order
language to a first-order language requires closure conversion, and no
proof of correctness with respect to denotational semantics covers
closure conversion for a Turing-complete source language.  Compiler
correctness requires that the denotation of a program be both
preserved and reflected by the compiler. Without reflection, a
compiler would be allowed to translate a diverging program into one
that halts.  For closure conversion, reflection is the hard
direction. We present the first denotational proofs of semantic
reflection for closure conversion applied to a Turing-complete
language, one proof with respect to a graph model and one with respect
to Scott domains. For the graph-model proof we develop a new logical
relation, the
\emph{approximant logical relation}, and for the Scott-domain proof we
adapt a logical relation of Reynolds.

As a first step in evaluating the tradeoffs, we compare mechanized
correctness proofs of an eight-pass compiler, from FPC, a typed
$\lambda$-calculus with recursive types, down to a typed low-level
intermediate language, in five semantic styles, each proof with
respect to one of the semantics. The five styles are the two
denotational semantics (graph model and Scott domain) and three
operational semantics (interpreter, reduction semantics, and abstract
machine). We find that no one semantics yields the shortest proofs for
every compiler pass, though the interpreter proof is the smallest
overall and the graph-model proof is the second smallest.  We identify
two factors that influence proof size. For operational semantics,
proof size tends to grow with how much the runs of the source and
target differ.  For denotational semantics, the proofs are succinct
when the correctness of a pass is an equation between meanings, as for
a source optimizer, and lengthy when a pass needs a logical relation,
as for the closure-conversion pass that separates each closure's code
from its environment.
